\section{The search tree and how it is pruned}
\label{sec:search}

\subsection{Canonical enumeration}

Summand permutations act on $\mathcal D_K$, so we fix a total order on $\Z^h$
(lexicographic) and represent each class by its sorted column tuple
$k_1\preceq k_2\preceq\cdots\preceq k_5$, repetitions allowed. Because
$\sum_ak_a=0$, the first four columns determine the fifth. The search is a
depth-first traversal of the tree whose depth-$m$ nodes are sorted prefixes
$(k_1,\dots,k_m)$, $m\le4$; the children of a node append a column
$k_{m+1}\succeq k_m$ from a precomputed pool, and a depth-4 node produces the
single leaf $k_5=-\sum_{a\le4}k_a$, kept only if it lies in the box and satisfies
$k_5\succeq k_4$. Every class of $\mathcal D_K$ is reached exactly once, which
gives a direct completeness argument; ambient-factor permutations are
\emph{not} quotiented during the search but only afterwards, when orbits are
compared (\cref{sec:validation}).

\begin{figure}[t]
\centering
\begin{tikzpicture}[
  node distance=0.55cm and 0.4cm,
  box/.style={draw, rounded corners=2pt, align=center, font=\small, inner sep=4pt, fill=#1},
  box/.default=white,
  rule/.style={font=\scriptsize, align=left, text width=6.6cm},
  arr/.style={-{Stealth[length=2.2mm]}, thick}]
\node[box=blue!6] (pool) {\textbf{Pool} $P_1\subset[-K,K]^h$\\ single-column rules};
\node[box=blue!6, below=of pool] (d1) {depth 1: $(k_1)$};
\node[box=blue!6, below=of d1] (d2) {depth 2: $(k_1,k_2)$};
\node[box=blue!6, below=of d2] (d3) {depth 3: $(k_1,k_2,k_3)$};
\node[box=blue!6, below=of d3] (d4) {depth 4: $(k_1,\dots,k_4)$};
\node[box=green!8, below=of d4] (leaf) {leaf: $k_5=-\sum_{a\le4}k_a$};
\node[box=orange!10, below=of leaf] (surv) {necessary-condition survivors};
\foreach \a/\b in {pool/d1,d1/d2,d2/d3,d3/d4,d4/leaf,leaf/surv} \draw[arr] (\a) -- (\b);
\node[rule, right=0.6cm of pool] {R1 individual slope feasibility (\cref{prop:percoef})\\
  R2 single-index range $-3|\Gam|\le\ind L\le0$\\ R3 pair-index domain (bitset)};
\node[rule, right=0.6cm of d2] {R4 component completion bound\\
  R5 partial single/pair index budgets\\
  R6 span slope obstruction (\cref{prop:monotone})};
\node[rule, right=0.6cm of d3] {R7 partial anomaly budget (\cref{prop:anomalybudget})\\ + R4--R6};
\node[rule, right=0.6cm of leaf] {R8 structure group \ref{c:structure}; R9 indices \ref{c:index}\\
  R10 anomaly \ref{c:anomaly}; R11 Bogomolov \ref{c:bogomolov}\\
  R6 complete slope test on the full span};
\node[rule, left=0.5cm of pool, text width=3.0cm, align=right] {before search:\\ divisibility $d_X\mid3|\Gam|$\\ (\cref{prop:divisibility})};
\end{tikzpicture}
\caption{The search tree. Each rule is a necessary condition (\cref{lem:soundness});
rules R4, R5 and R6 are monotone and prune entire subtrees. Rules R3 and R4 are
enforced by construction through bitset-indexed candidate lists rather than
tested inside the loop.}
\label{fig:tree}
\end{figure}

\subsection{The rules}

\Cref{tab:rules} lists every rule, its status and where it acts. We give the
proofs of those that are not immediate.

\begin{table}[t]
\centering\small
\begin{tabularx}{\textwidth}{@{}lXll@{}}
\toprule
& Rule & Status & Acts at \\
\midrule
R1 & $\mu_k$ has a zero in $T_+$ (signed-matrix test; exact per \cref{prop:percoef}) & exact (bench.) & pool \\
R2 & $-3|\Gam|\le\ind\Oc_X(k)\le0$ (\st{GUT\_FULL}) & necessary & pool \\
R3 & declared domain: $\ind(L_a\otimes L_b)\le0$ for all pairs & domain restriction & pool, prefix \\
R4 & $|\sum_{a\le m}k_a^i|\le(5-m)K$ & necessary & prefix \\
R5 & partial sums of single and pair indices lie in $[-3|\Gam|,0]$ & necessary & prefix \\
R6 & $\spn_\R\{k_a\}\cap P=\{0\}$ (complete slope obstruction) & necessary & prefix, leaf \\
R7 & partial anomaly budget & necessary & prefix \\
R8 & no proper zero-sum subset \ref{c:structure} & exact & leaf \\
R9 & $\ind V=-3|\Gam|$ ($\ind\wedge^2V$ then automatic, \cref{prop:indexidentity}) & exact & leaf \\
R10 & $c_{2,r}-c_2(V)_r\ge0$ & necessary & leaf \\
R11 & Bogomolov pairing & necessary & leaf \\
R12 & $d_X\mid 3|\Gam|$ & exact, unbounded & before search \\
\bottomrule
\end{tabularx}
\caption{Pruning rules. R3 restricts the declared domain to the historical
domain of \cite{ACGLP2013}; it is not claimed necessary, and the independent
replays (\cref{sec:validation}) also run with it removed. R1 is exact on
geometries with $d_{ijj}=0$, which includes both benchmarks.}
\label{tab:rules}
\end{table}

\paragraph{R2 (single-index range).} At a common zero of the slopes, a
nontrivial line bundle $L_a$ has $h^0=h^3=0$, since a line bundle of slope zero
with a nonzero section is trivial, and \ref{c:structure} excludes trivial
summands. No anti-families means $h^2(X,V)=0$, hence $h^2(L_a)=0$ and
$\ind L_a=-h^1(L_a)\le0$. The five indices sum to $-3|\Gam|$, so each lies in
$[-3|\Gam|,0]$, and so does every partial sum (R5).

\paragraph{R4 (component completion).} With $m$ columns fixed, the remaining
$5-m$ columns have entries in $[-K,K]$ and must cancel the partial sum, so
$|\sum_{a\le m}k_a^i|\le(5-m)K$ by the triangle inequality.

\subsection{The slope obstruction}
\label{sec:slope}

Poly-stability of a sum of line bundles with $c_1(V)=0$ means a common zero of
the five slopes:
\begin{equation}\tag{SZ}\label{eq:SZ}
\exists\,t\in T_+:\qquad \mu_{k_a}(t)=t^\top M(k_a)\,t=0\qquad(a=1,\dots,5).
\end{equation}

\begin{lemma}[Sign obstruction]\label{lem:sign}
If $M\in\Sym_h(\Z)$ is nonzero and entrywise nonnegative (or nonpositive), then
$t^\top Mt\neq0$ for every $t\in T_+$.
\end{lemma}
\begin{proof}
Each term $M_{jl}t^jt^l$ has the same weak sign and at least one is strictly
nonzero, because $t^jt^l>0$.
\end{proof}

\begin{proposition}[Combined test]\label{prop:combined}
If \eqref{eq:SZ} holds then for every $c\in\R^5$ the matrix
$M(\sum_ac_ak_a)=\sum_ac_aM(k_a)$ is zero or has entries of both signs.
Contrapositively, a single $c\in\Z^5$ for which $M(\sum_ac_ak_a)$ is nonzero and
one-signed proves that $V$ is not poly-stable anywhere in $T_+$.
\end{proposition}
\begin{proof}
At the witness $t_*$, $t_*^\top M(\sum_ac_ak_a)t_*=\sum_ac_a\mu_{k_a}(t_*)=0$;
apply \cref{lem:sign}.
\end{proof}

The scan of \cite{ACGLP2013} applies \cref{prop:combined} with
$c\in\{-1,0,1\}^5$. This is \emph{sound} but not \emph{complete}: passing the
test for all $c$ in a finite box does not imply \eqref{eq:SZ}. (Since
$\sum_ak_a=0$, $M(\sum_ac_ak_a)$ depends only on the class of $c$ modulo
$\Z(1,\dots,1)$; the 121 vectors of the box with first nonzero entry positive
fall into 106 classes, and the box is exactly the set of classes of spread
$\max c-\min c\le2$.) The quantified statement, however, is decidable exactly.

\begin{proposition}[Exact obstruction cone]\label{prop:cone}
Let $S=\spn_\R\{k_1,\dots,k_5\}\subseteq\R^h$ and
$P=\{k\in\R^h: M(k)_{jl}\ge0\ \forall j,l\}$, a rational polyhedral cone. Some
$c\in\Z^5$ yields a one-signed nonzero combined matrix if and only if $S\cap P$
contains a point $k$ with $M(k)\neq0$.
\end{proposition}
\begin{proof}
$M$ is linear, so $\{M(\sum_ac_ak_a):c\in\R^5\}=M(S)$, and $P$ is the preimage of
the entrywise-nonnegative matrices; $-c$ handles the nonpositive case. A
rational point of the rational cone $S\cap P$ can be scaled to an integer point
$k$, and $c$ is recovered by solving the consistent rational system
$\sum_ac_ak_a=k$, then clearing denominators.
\end{proof}

On geometries with $d_{ijj}=0$ for all $i,j$---both benchmarks---$M(k)$ has zero
diagonal, and nothing is lost per coefficient vector:

\begin{proposition}[Per-coefficient optimality]\label{prop:percoef}
If $d_{ijj}=0$ for all $i,j$, then for any $k\in\Z^h$, the form $\mu_k$ has a
zero in $T_+$ if and only if $M(k)$ is zero or has entries of both signs.
\end{proposition}
\begin{proof}
($\Rightarrow$) is \cref{lem:sign}. ($\Leftarrow$): the diagonal vanishes, so
suppose $M_{jl}>0$ and $M_{j'l'}<0$ at off-diagonal positions. Let $t(T)$ equal
$T$ at positions $j,l$ and $1$ elsewhere; then
$\mu_k(t(T))=2M_{jl}T^2+O(T)\to+\infty$. The analogous path at $(j',l')$ gives
$-\infty$. $T_+$ is connected and $\mu_k$ continuous, so $\mu_k$ vanishes in $T_+$.
\end{proof}

In particular the single-bundle condition ``$\mu(L)=0$ somewhere in the
interior of the Kähler cone'' is decided exactly on these geometries by the
signed-matrix test (R1), and the only incompleteness of the finite test is the
truncation of the $c$-search, which \cref{prop:cone} removes.

\begin{proposition}[Monotonicity]\label{prop:monotone}
If $S'\subseteq S$ and $S'\cap P$ contains a point with $M\ne0$, so does $S\cap P$.
Hence the test applied to the span of any prefix $(k_1,\dots,k_m)$ prunes every
completion of that prefix.
\end{proposition}
\begin{proof}
Immediate from $S'\cap P\subseteq S\cap P$.
\end{proof}

To decide $S\cap P\neq\{0\}$ exactly we need $P$ to be pointed and its extreme
rays.

\begin{lemma}[The tetraquadric cone]\label{lem:rays}
On $\Xa$, $M(k)_{jl}=2(k^m+k^n)$ for $j\ne l$, where $\{m,n\}$ is the
complement of $\{j,l\}$ in $\{1,2,3,4\}$, and $M(k)_{jj}=0$. Hence
$P=\{k\in\R^4:k^m+k^n\ge0\ \forall m<n\}$, $M(k)=0$ only for $k=0$, and $P$ is
pointed with exactly eight extreme rays: the coordinate vectors $e_i$ and the
four vectors with one entry $-1$ and the others $+1$.
\end{lemma}
\begin{proof}
The formula follows from $d_{ijl}=2$ for distinct indices. If all pairwise
sums vanish then $k^1+k^2=k^1+k^3=k^2+k^3=0$ forces $k^1=k^2=k^3=0$ and then
$k^4=0$, so $P\cap(-P)=\{0\}$. A ray of $P$ is extreme iff the constraints tight
on it have rank three. For $e_1$ the tight constraints are the three pairs not
containing $1$; for $(-1,1,1,1)$ the three pairs containing $1$; each set has
rank three. Enumerating the $\binom63=20$ triples of constraints shows there are
no others.
\end{proof}

For $\Xb$ the analogous computation (exact, by enumeration) gives $2h=10$
extreme rays: the $e_i$ and the vectors with one entry $-2=-(h-3)$ and the
others $+1$. In general the implementation enumerates extreme rays of $S\cap P$
exactly in integer arithmetic and declines to prune if the pointedness check
fails.

The primal enumeration costs a number of kernel computations that grows with
the dimension of $S$ and dominated the running time on $\Xb$ (90\% of a profiled
scan). The dual form is much cheaper exactly where the search spends its time.

\begin{proposition}[Dual decision]\label{prop:dual}
Let $P$ be pointed with extreme rays $r_1,\dots,r_s$, let $A$ be a matrix whose
rows span $S^\perp$, and $d=\dim S^\perp=h-\rk S$. Then
\[
S\cap P=\{0\}\iff\exists\,\varphi\in S^\perp:\ \langle\varphi,r_i\rangle>0\ \forall i
\iff 0\notin\conv\{Ar_1,\dots,Ar_s\}\subset\R^d .
\]
\end{proposition}
\begin{proof}
As $P$ is pointed its base $B=\conv\{r_i\}$ is compact with $0\notin B$. If
$S\cap P=\{0\}$ then $S\cap B=\emptyset$, and the compact convex set $B$ is
strictly separated from the closed subspace $S$ by a functional that must vanish
on $S$, i.e.\ lies in $S^\perp$, and is positive on $B$. Conversely such a
$\varphi$ vanishes on $S$ and is positive on $P\setminus\{0\}$. The last
equivalence is Gordan's alternative~\cite{Schrijver1986} for the projected
points $Ar_i$.
\end{proof}

The dimension $d$ \emph{shrinks} as the prefix grows. At the deepest and most
numerous nodes $d\in\{1,2\}$: for $d=1$ the test is whether the projected
scalars take both signs ($O(s)$), for $d=2$ an angular-gap test
($O(s\log s)$). The primal enumeration remains the reference implementation,
the fallback for $d>2$, and the producer of the explicit integer witness $k_c$
recorded in every rejection certificate. The two were differentially tested on
12,000 random column sets across every span rank with no disagreement.

\begin{example}[A spurious survivor of the finite test]\label{ex:spurious}
On $\Xa$ with $\Gam=\Z_2$ at $K=3$, the charge matrix with columns
\[
(-3,-1,1,2),\ (0,1,-1,0),\ (0,1,-1,0),\ (1,1,-1,-1),\ (2,-2,2,-1)
\]
satisfies every other declared condition and passes the test for all
$c\in\{-1,0,1\}^5$; a numerical optimiser stalls at slope residual
$3.6\times10^{-7}$. The complete test rejects it with the integer witness
$k_c=(1,0,0,0)\in S$, for which $M(k_c)$ is entrywise nonnegative and nonzero.
It is absent from the published catalog. Its orbit is the only difference
between the finite and the complete test at $K=3$ (8 versus 7 orbits) and at
$K=4$ (11 versus 10), and removing it is what gives exact catalog agreement
(\cref{sec:validation}).
\end{example}

Soundness was tested by a counterexample search: 12,000 randomly generated
bundles carrying an exact common slope zero on $\Xa$ and $\Xb$ were never
rejected; an independently written Fourier--Motzkin implementation agreed with
the extreme-ray implementation on 3,000 random column sets; and every published
model and every survivor with an exact slope witness is retained.

\subsection{Beyond sign tests: a Pólya hierarchy and an exact decision}
\label{sec:polya}

The sign test is exact per combination only when $d_{ijj}=0$
(\cref{prop:percoef}), and even then a bundle can pass every sign test while its
five slopes have no \emph{simultaneous} zero. Two further exact tools address
this; both produce certificates checked in exact arithmetic.

\begin{proposition}[Pólya certificates]\label{prop:polya}
If for some $c$ and some $N\ge0$ the form
$(t_1+\dots+t_h)^N\sum_ac_a\mu_{k_a}(t)$ has all coefficients nonnegative and not
all zero, then \eqref{eq:SZ} has no solution. Every form strictly positive on the
closed simplex admits such an $N$; $N=0$ is the sign test.
\end{proposition}
\begin{proof}
A nonzero polynomial with nonnegative coefficients is positive on the open
orthant, and $(\sum_it_i)^N>0$ there; at a common zero the combination would
vanish. The converse direction is Pólya's theorem.
\end{proof}

The coefficients are linear in $k_c=\sum_ac_ak_a$, so each level of the hierarchy
is again a polyhedral cone and a certificate is found by linear programming and
verified exactly. On geometries with $d_{ijj}=0$ the hierarchy provably adds
nothing beyond $N=0$ (by \cref{prop:percoef}); on the other catalog geometries we
tested it added no certificate for charges in $[-2,2]$ up to $N=4$. The
obstruction for the bundles below, where it exists, is therefore \emph{joint}.

\begin{proposition}[Exact decision in rank $h-1$]\label{prop:decision}
Suppose $\rk K=h-1$. Set $t_h=1$, which loses no interior point, and compute a
lexicographic Gröbner basis of $h-1$ independent slopes. If it is in shape
position, $t_i=g_i(x)$ and $p(x)=0$; isolate every real root of every irreducible
factor of $p$ in an exact rational interval and enclose the other coordinates by
rigorous interval arithmetic. A root at which all coordinates are positive, and
modulo whose minimal polynomial every slope vanishes identically, is an exact
real-algebraic witness of poly-stability. If no real root is interior, \eqref{eq:SZ}
has no solution.
\end{proposition}
\begin{proof}
In rank $h-1$ the common zeros with $t_h=1$ are the solutions of the $h-1$
independent equations, which in shape position are finitely many and given
exactly by the roots of $p$; interval enclosures that exclude zero decide each
coordinate's sign rigorously.
\end{proof}

When the basis is not in shape position the decision proceeds factor by factor:
factors of $p$ without a positive real root are discarded (an interior point has
$t_{h-1}>0$); a rational root is substituted and the smaller system decided
recursively; over an irrational factor $f$ the remaining coordinates are solved
linearly in $\Q[x]/(f)$ when possible. Anything else, including positive-dimensional
zero sets (rank below $h-1$), is reported as \st{UNDECIDED}. On the tetraquadric the procedure certifies the published model
with witness $t=(1,1,1,1)$, and the spurious candidate of \cref{ex:spurious}, whose
zero set is positive-dimensional, is excluded by the sign test instead: the tools
are complementary.

\subsection{Index identity and divisibility}

\begin{proposition}[Index identity]\label{prop:indexidentity}
On any Calabi--Yau threefold, if $V=\bigoplus_{a=1}^5L_a$ has $c_1(V)=0$ then
$\ind\wedge^2V=\ind V$.
\end{proposition}
\begin{proof}
The linear terms vanish: $\sum_ac_2\cdot k_a=0$ and each $a$ occurs in four
pairs. For the cubic terms,
$\sum_{a<b}D(k_a+k_b)^3=4\sum_aD(k_a)^3+3\sum_{a\ne b}D(k_a,k_a,k_b)$ and
$\sum_{b\ne a}k_b=-k_a$, so the sum is $4\sum_aD(k_a)^3-3\sum_aD(k_a)^3$.
\end{proof}
This is stated in \cite{ACGLP2013} after their equation~(4.7); we use it to
drop a redundant filter and keep it as a free consistency check. (For rank $n$
the coefficient is $n-4$.)

\begin{proposition}[Divisibility obstruction]\label{prop:divisibility}
Let $d_X=\gcd\{\ind\Oc_X(k):k\in\Z^h\}$. If $d_X\nmid3|\Gam|$, no sum of line
bundles on $X$ has $\ind V=-3|\Gam|$, for any charges. Moreover
$d_X=\gcd\{\ind\Oc_X(k):k\in\{0,1,2,3\}^h\}$.
\end{proposition}
\begin{proof}
$\ind V=\sum_a\ind L_a\in d_X\Z$. For the second claim, $\ind$ is an
integer-valued polynomial of degree three, hence a $\Z$-combination of
$\prod_i\binom{k_i}{\alpha_i}$ with $|\alpha|\le3$; its values generate the same
subgroup as the coefficients, which are integer combinations of finite
differences at points of $\{0,1,2,3\}^h$.
\end{proof}

Over the 68 favourable CICYs, $d_X\in\{1,2,3,5\}$ occurs $19,44,4,1$ times.
Combined with the symmetry orders recorded for each geometry, the obstruction
predicts \emph{zero} models for exactly two pairs, $(7761,|\Gam|=2)$ with $d_X=5$
and $(7447,|\Gam|=5)$ with $d_X=2$; the published catalog has zero models for
both, and no catalog entry with models contradicts it. The prediction is not
vacuous: geometries with odd $|\Gam|$ and $d_X\mid3|\Gam|$ do carry models.
The rule is elementary; we have not found it stated in the sources we examined.
It is the only rule here whose conclusion carries no bound on the charges.

\subsection{The partial anomaly budget}

\begin{proposition}\label{prop:anomalybudget}
Assume \ref{c:c1}, a trivial hidden bundle and nef $J_r$. Then \ref{c:anomaly}
implies the separable inequality
\begin{equation}\label{eq:anomaly-separable}
\textstyle\sum_{a=1}^5q_r(k_a)\ \ge\ -2c_{2,r}\qquad\text{for all }r .
\end{equation}
Let $P_1$ be a pool containing every column any valid completion can use, and
$Q_r=\max_{k\in P_1}q_r(k)$. If for a prefix $(k_1,\dots,k_m)$ and some $r$
\[
\textstyle\sum_{a\le m}q_r(k_a)+(5-m)Q_r<-2c_{2,r},
\]
no completion satisfies \eqref{eq:anomaly-separable} and the branch is pruned.
\end{proposition}
\begin{proof}
With $\sum_ak_a=0$ the Whitney formula gives
$c_2(V)_r=-\tfrac12\sum_aq_r(k_a)$, so $c_{2,r}-c_2(V)_r\ge0$ is
\eqref{eq:anomaly-separable}. Each remaining column contributes at most $Q_r$.
\end{proof}

The pool in the implementation is the retained single-column list, itself
defined by necessary conditions, and the derived fifth column is required to
lie in it. The bound is inactive for $K\le2$, since small charges cannot reach
it, and first bites at $K=3$ (\cref{tab:ablation}).

\subsection{Implementation}

Two features matter for the complexity. First, the retained pool is indexed by
bitsets: for each component $i$ and threshold, the set of columns whose
component lies above or below it, and for each column the set of compatible
partners under R3. Rules R3 and R4 are then enforced \emph{by construction}---a
node iterates only over the intersection of the relevant bitsets---rather than
tested inside the loop. This cut fourth-column attempts by factors of 19--28
and wall time by $3$--$7.5\times$ at identical survivor sets, which was
verified by re-running every recorded domain and requiring exact set equality.
Second, the span of the prefix is maintained incrementally in a fraction-free
integer row-reduced form, which is also the memoisation key for the slope
decision, so repeated spans are decided once. Every run is checkpointed; a
checkpoint embeds the SHA-256 of the source and inputs and refuses to resume
under changed code.

\subsection{Measured effect of the pruning}

\Cref{tab:funnel} shows the funnel at $K=2$ and \cref{tab:ablation} the effect
of each rule and of $K$.

\begin{table}[t]
\centering\small
\begin{tabular}{@{}lr@{}}
\toprule
stage ($\Xa$, $\Gam=\Z_2$, $K=2$, finite slope test, pre-indexing engine) & count \\
\midrule
ordered four-column assignments (nominal) & 152,587,890,625 \\
individual columns in the box & 625 \\
retained single columns (R1, R2) & 235 \\
retained depth-2 prefixes & 10,573 \\
retained depth-3 prefixes & 123,629 \\
fourth-column attempts & 7,875,639 \\
retained depth-4 prefixes & 252,154 \\
complete canonical bundles & 38,161 \\
\quad no proper zero-sum subset (R8) & 24,043 \\
\quad $\ind V=-6$ (R9) & 6,261 \\
\quad all pair indices in range (R3) & 1,356 \\
\quad anomaly (R10) & 1,260 \\
\quad slope (R6) & \textbf{120} \\
\bottomrule
\end{tabular}
\caption{Funnel of the search. Prefix counters count prefixes, not full states,
and must not be added as if they partitioned the nominal domain. The ratio of
nominal assignments to materialised bundles, about $4\times10^6$, includes
canonicalisation and is not a pure physics pruning factor; at the last level the
slope obstruction alone removes a factor $10.5$.}
\label{tab:funnel}
\end{table}

\begin{table}[t]
\centering\small
\begin{tabular}{@{}rlrrrrrr@{}}
\toprule
$K$ & slope rule & wall (s) & RSS (MiB) & 4th-col.\ attempts & bundles built & survivors & orbits \\
\midrule
1 & finite & 0.05 & 25.4 & 4,636 & 713 & 0 & 0 \\
1 & complete & 0.09 & 25.3 & 2,535 & 334 & 0 & 0 \\
2 & finite & 4.00 & 33.2 & 420,266 & 38,161 & 120 & 5 \\
2 & complete & 3.99 & 37.2 & 170,121 & 10,097 & 120 & 5 \\
3 & finite & 15.88 & 50.5 & 2,304,135 & 146,425 & 192 & \textbf{8} \\
3 & complete & 18.42 & 68.3 & 963,549 & 52,800 & 168 & 7 \\
4 & finite & 36.65 & 65.6 & 7,147,830 & 294,130 & 264 & \textbf{11} \\
4 & complete & 43.62 & 105.6 & 3,236,246 & 118,577 & 240 & 10 \\
5 & complete & 85.44 & 141.0 & 6,537,468 & 162,761 & 240 & 10 \\
6 & complete & 118.88 & 193.8 & 10,199,673 & 206,847 & 240 & 10 \\
7 & complete & 179.15 & 188.9 & 14,516,351 & 231,946 & 240 & 10 \\
\bottomrule
\end{tabular}
\caption{$\Xa$, $\Gam=\Z_2$, target \st{GUT\_FULL}, bitset-indexed engine.
Orbits are ambient-row orbits of the column-canonical survivors. The finite test
retains one spurious orbit at $K=3$ and $K=4$ (\cref{ex:spurious}). With the dual
decision of \cref{prop:dual} the complete rule at $K=4$ takes 28.8\,s, faster
than the finite test. The partial anomaly budget (\cref{prop:anomalybudget})
first acts at $K=3$, pruning 3,734 depth-3 and 166,040 depth-4 branches under the
complete rule and reducing built bundles from 75,534 to 52,800.}
\label{tab:ablation}
\end{table}

Two observations. The complete slope rule, applied at every depth, removes
2,730 of the 10,573 depth-2 prefixes already at $K=2$, and cuts the bundles
built by factors of $2.8$ at $K=3$ and $2.5$ at $K=4$. And the pruning improves
with $K$ rather than degrading: from $K=1$ to $3$ the nominal domain grows by
factors $3,544$ and $218$ while the bundles built grow by $30$ and $5.2$.
